% !TeX root = ../概率论与数理统计笔记.tex
% 由主文档通过 \input 载入；请编译主文档。
\setcounter{chapter}{9} % 保留原书章号。
\chapter{总体、样本、统计量与抽样分布}
\label{chap:ps-sampling}

本章从已知概率模型转向未知总体的学习问题，为所有推断方法建立共同框架。
先修内容为分布、数字特征与极限定理；学习目标是区分总体量、样本统计量及其抽样分布。

\section{总体、统计模型与抽样机制}
\label{sec:ps-statistical-model}
定义\emph{总体}（population）、\emph{统计模型}（statistical model）$\{\mathbb P_\theta:\theta\in\Theta\}$ 与参数空间。
本书将来自同一分布的独立同分布样本称为\emph{简单随机样本}（independent and identically distributed sample），并说明有限总体无放回简单随机抽样是另一种机制，其观测通常不独立。
区分随机抽样、随机分组与便利取样，说明推断范围受数据获得方式约束。

\section{统计量与经验分布}
\label{sec:ps-statistics-empirical}
定义不含未知参数的\emph{统计量}（statistic），比较统计量、估计量和估计值。
介绍\emph{经验分布函数}（empirical distribution function），以及直方图、箱线图和分位数等描述方式。
说明描述数据的图表与关于总体的推断有不同任务，分组和汇总可能掩盖数据结构。

\section{样本均值、样本方差与样本矩}
\label{sec:ps-sample-moments}
对 $n\geq2$ 的样本定义样本均值 $\overline X_n$ 与样本方差
\begin{equation}
  S^2=\frac{1}{n-1}\sum_{j=1}^n(X_j-\overline X_n)^2.
  \label{eq:ps-sample-variance}
\end{equation}
将估计量的期望等于目标总体量称为\emph{无偏性}（unbiasedness）。
在独立同分布且二阶矩有限时，证明样本均值和样本方差分别无偏估计总体均值和总体方差，解释分母 $n-1$ 的来源。
比较样本方差、样本二阶中心矩和样本均值的方差，避免把它们当作同一对象。

\section{正态总体的抽样分布}
\label{sec:ps-normal-sampling}
定义卡方分布（chi-square）、$t$ 分布（Student's $t$）与 $F$ 分布，并说明构造中标准正态变量或卡方变量的独立性要求。
对正态独立同分布样本建立均值与样本方差的分布及独立性；推导单样本 $t$ 统计量和独立正态样本的方差比。
自由度的含义通过独立信息与线性约束解释，相关正交分解证明按所需线性代数程度展开。

\section{次序统计量与样本分位数}
\label{sec:ps-order-statistics}
定义\emph{次序统计量}（order statistic）$X_{(1)}\leq\cdots\leq X_{(n)}$。
对独立同分布且具有密度的样本，推导第 $k$ 个次序统计量的密度，并说明 $1\leq k\leq n$。
研究极差、中位数和样本分位数，区分离散并列值与连续无并列值情形。

\section{例题、习题与学习检查}
规划由正态样本构造 $t$ 统计量、求均匀样本最大值分布、比较有偏与无偏样本方差公式。
习题安排“是否为统计量”的辨析、抽样分布推导与抽样偏差讨论。
学习检查要求说明重复抽样中变化的对象，并区分样本标准差和估计量的标准误。






